\documentclass[Screen16to9,17pt]{foils}
\usepackage{zencurity-slides}

\selectlanguage{english}

\usepackage{tabularx}

% Reused from simulated-ddos-workshop + pentest-I-introduction (obj 9 + resilience).

\begin{document}

\mytitlepage{Network Attacks II -- Hands-on \& Defence}
{Cyber Akademiet, 2026}

\slide{Goals for today}
\begin{list2}
\item Demonstrate network-attack techniques in an authorised lab
\item Generate floods with hping3 and craft packets with Scapy
\item Design defences: filtering, rate/connection limits, monitoring, CDN/DDoS protection
\end{list2}

\slide{Plan for today}
\begin{list2}
\item Lab setup, legal issues and scope
\item Packet generators, hping3 and Scapy -- hands-on
\item Designing the protection: filtering and rate limits
\item Results: before and after filtering
\end{list2}

\slide{Time schedule}
\begin{list2}
\item 13:15--14:45 Lab, hping3, Scapy -- \textbf{hands-on} (booklet)
\item 14:45--15:00 Break
\item 15:00--16:00 Defences: filtering, rate limits, results
\item 16:00--16:15 Wrap-up
\end{list2}

\slide{Kali Linux the new backtrack}

\hlkimage{16cm}{kali-linux.png}

\begin{list1}
\item Kali \link{http://www.kali.org/}
\item Use the documentation on Kali website for installation and management of the operating system
\item Use the web sites belonging to projects for documentation about tools, \\
like \link{https://nmap.org/} and \link{http://www.hping.org/}
\end{list1}

\slide{Testing network the legal issues}

{\bfseries Straffelovens paragraf 263 Stk. 2. Med bøde eller fængsel indtil 1 år og 6 måneder straffes den, der uberettiget skaffer sig adgang til en andens oplysninger eller programmer, der er bestemt til at bruges i et informationssystem. }

\begin{list2}
\item Danish law about hacking
\item Please check with your legal department, or be careful
\item We {\bf always} contact network between us and the network to be tested
\item Be good netizens
\end{list2}

\slide{Hackerlab setup}

\hlkimage{5cm}{hacklab-1.png}

\begin{list2}
\item I recommend getting a hackerlab running on your laptop
\item Hardware: modern laptop which has CPU virtualization\\
Dont forget to check BIOS settings for virtualization
\item Software: your favorite OS: Windows, Mac, Linux
\item Virtualization software: VMware, Virtual box, HyperV
\item Hacker software: Kali as a Virtual Machine \link{https://www.kali.org/}
\item Tools -- we will use two packet generators today hping3 and t50
\item Testing production systems would probably be better with physical host\\
but only a small PC is needed even for 10Gbps testing
\end{list2}

\slide{Scanning and Attacking -- \bf Pressure Points and Scope}

% Drawing showing "network"

\hlkimage{21cm}{generic-network-pressure-points.pdf}

\begin{list2}
\item In scope for me is everything that could adversely affect the network
\item Typical scope IPv4: Link network /28 or /26 and a hosting network /24 or even /22
\item Typical scope IPv6: Link network /64 (bad) or /127 (RFC9099) and a hosting network /48 with subnets
%\item Domain Name Service outside of the network is ofc also attack vector, ignored today
\end{list2}

\slide{Prepare for the testing}


\hlkimage{10cm}{trinity-nmapscreen-hd-cropscale-418x250.jpg}

\begin{list2}
\item Portscan the whole linknet and hosting range in IPv4
\item Ask about IPv6 ranges in use, specific subnets and IPv6 addresses\\
We can guess from traceroute, Nmap test first 100 addresses in each subnet etc. but easier to ask
\item Portscan the whole linknet - identify provider devices, hosting network devices, type of device router/firewall
\item Also Thank you Fyodor (Gordon Lyon) and contributors for Nmap!
\end{list2}

\slide{Detailed Scope and Plan}

Select a few targets for monitoring and attacks, from the port scans

Best case would be to have:
\begin{list2}
\item Ping ICMP allowed to a provider router and hosting firewall -- is the connection alive
\item TCP port with service checks, HTTP being attacked and one not being attacked
\item UDP port with service check, DNS is a favourite -- ask for localhost/127.0.0.1
\item Put monitoring on these, a week before testing is nice
\item Agree on a day or night for testing, inform participants and system owners
\end{list2}

\slide{Performing the DDoS test }

So lets break this task down into:
\begin{list2}
\item Use Nmap to port scan the network
\item Run testing tool \verb+hping3+, \verb+t50+ or \verb+penguinping+
\item Go through all scenarios before making changes to environment
\item Expect things to break, investigate, repeat failed scenarios
\vskip 1cm
\item BTW we usually schedule this for night time! There WILL be interruptions
\vskip 1cm
\item How do you run the tools, since they all have so many options?\\
\verb+man nmap | enscript -o test.ps+ result in 54 pages,
\verb+hping3+ has about 80 options and \verb+t50+ has even more
\end{list2}

% Slides from courses about Hping3

\slide{Packet processing in firewalls -- detailed view}

% Picture with pipeline from Juniper SRX
\hlkimage{23cm}{srx-firewall-flow-based-processing.png}
\emph{Traffic Processing on SRX Series Devices Overview}\\ {\scriptsize
\link{https://www.juniper.net/documentation/us/en/software/junos/flow-packet-processing/topics/topic-map/security-srx-devices-processing-overview.html}}

\slide{Packet generators}

\hlkimage{3cm}{penguinping.jpg}


\begin{list2}
\item Scapy is an example packet generator, allowing you to use Python
\item It is very flexible and often \emph{fast enough}, but other times, not fast enough
\item Lots of other tools exist, some work as kernel modules even
\item I love hping3 and t50 -- and use them for stress testing firewalls and devices
\item Started building a replacement or \verb+penguinping+
\end{list2}
\centerline{\Large $\heartsuit$ love internet packets $\heartsuit$}

\slide{hping3 packet generator}

\begin{alltt}\footnotesize
usage: hping3 host [options]
  -i  --interval  wait (uX for X microseconds, for example -i u1000)
      --fast      alias for -i u10000 (10 packets for second)
      --faster    alias for -i u1000 (100 packets for second)
      --flood      sent packets as fast as possible. Don't show replies.
UDP/TCP
  -s  --baseport   base source port             (default random)
  -p  --destport   [+][+]<port> destination port(default 0) ctrl+z inc/dec
  -L  --setack     set TCP ack
  -F  --fin        set FIN flag
  -S  --syn        set SYN flag
...
\end{alltt}

\begin{list2}
\item Hping3 packet generator is a very flexible tool to produce simulated DDoS traffic with specific charateristics
\item Home page: \link{http://www.hping.org/hping3.html} Source repository \link{https://github.com/antirez/hping}
\item My fork with IPv6 and VXLAN branches added \link{https://github.com/kramse/hping-2018}
\item Used to be my primary DDoS testing tool, now I am moving to PenguinPing / MoonGen
\end{list2}

\slide{t50 packet generator}


\begin{alltt}\footnotesize
# {\bf t50 -?}
T50 Experimental Mixed Packet Injector Tool 5.4.1
Originally created by Nelson Brito <nbrito@sekure.org>
Maintained by Fernando Mercês <fernando@mentebinaria.com.br>

Usage: T50 <host> [/CIDR] [options]

Common Options:
    --threshold NUM        Threshold of packets to send     (default 1000)
    --flood                This option supersedes the 'threshold'
...
6. Running T50 with '--protocol T50' option, sends ALL protocols sequentially.
# {\bf t50 -? | wc -l}
264
\end{alltt}

\begin{list2}
\item T50 packet generator, another high speed packet generator
can easily overload most firewalls by producing randomized traffic with multiple protocols like IPsec, GRE, MIX \\
home page: \link{http://t50.sourceforge.net/resources.html}
\item Very fast and breaks most firewalls when flooding all protocols, easy 800k pps/400Mbps
\end{list2}

\slide{Running Attacks with hping3}

\begin{alltt}\small
# export CUST_IP=192.0.2.138
# date;time hping3 -q -c 1000000  -i u60 -S -p 80 $CUST_IP &
Thu Jan 21 22:37:06 CET 2022
HPING 192.0.2.1 (eth0 192.0.2.1): S set, 40 headers + 0 data bytes

--- 192.0.2.1 hping statistic ---
1000000 packets transmitted, 999996 packets received, 1% packet loss
round-trip min/avg/max = 0.9/7.0/1005.5 ms

real    1m7.438s
user    0m1.200s
sys     0m5.444s
\end{alltt}

\vskip 1cm
\centerline{Dont forget to do a killall hping3 when done \smiley }

\slide{Test-cases / Scenarios}

The minimal run contains at least these:
\begin{list2}
\item SYN flood: \verb+hping3 -q -c 1000000  -i u60 -S -p 80 $CUST_IP &+
\item SYN+ACK: \verb+hping3 -q -c 1000000  -i u60 -S -A -p 80 $CUST_IP &+
\item ICMP flood: \verb+hping3 -q -c 1000000 --flood -1 $CUST_IP &+
\item UDP flood: \verb+hping3 -q -c 1000000 --flood -2 $CUST_IP &+
\item Near end of test we also throw in the joker to {\bf kill firewalls}\\
\verb+t50 --flood --protocol T50 $CUST_IP+
\end{list2}

While testing I use the tool \verb+ifpps+ tool from the Netsniff-ng package \link{http://netsniff-ng.org/} to monitor sending speed, or you can use your router/switch -- Junos \verb+monitor interface+

\slide{Test-cases / Scenarios, continued Spoof Source}

Spoofed packets src=dst=target \smiley

Flooding with spoofed packet source, within hosting range, \\
may result in your {\bf single packet} returning answer to another inside network + ARP traffic

\begin{alltt}\small

-a --spoof hostname
    Use this option in order to set a fake IP  source  address,  this
    option ensures that target will not gain your real address.
\end{alltt}

\verb+hping3 -q --flood -p 80 -S -a $CUST_IP $CUST_IP+

Preferably using a test-case you know fails, to see effect

Still amazed how often this works

\slide{Test-cases / Scenarios, continued Small Fragments}

Using the built-in option -f for hping

\begin{alltt}\small
-f --frag
    Split  packets  in more fragments, this may be useful in order to test IP
    stacks fragmentation performance and to test if some packet filter is  so
    weak  that  can  be  passed using tiny fragments (anachronistic). Default
    {\bf 'virtual mtu' is 16 bytes}. see also --mtu option.
\end{alltt}

\begin{list1}
\item \verb+hping3 -q --flood -p 80 -S -f $CUST_IP+
\item Similar process with bad checksum and Bad fragment offset
\end{list1}

\slide{Snippets of Scapy}

First create VXLAN header and inside packet
\begin{minted}[fontsize=\small]{python}
vxlanport=4789     # RFC 7384 port 4789, Linux kernel default 8472
vni=37             # Usually VNI == destination VLAN
vxlan=Ether(dst=routermac)/IP(src=vtepsrc,dst=vtepdst)/
   UDP(sport=vxlanport,dport=vxlanport)/VXLAN(vni=vni,flags="Instance")
broadcastmac="ff:ff:ff:ff:ff:ff"
randommac="00:51:52:01:02:03"
attacker="185.27.115.666"
destination="10.0.0.10"
# port is the one we want to contact inside the firewall
insideport=53
testport=54040
packet=vxlan/Ether(dst=broadcastmac,src=randommac)/IP(src=attacker,
    dst=destination)/UDP(sport=testport,dport=insideport)/
    DNS(rd=1,id=0xdead,qd=DNSQR(qname="www.wikipedia.org"))
\end{minted}

{\footnotesize Fun fact, Unbound on OpenBSD reply to DNS requests received in Ethernet packets with broadcast destination and IP destination being the IP of the server}

\slide{Send and receive - from another source}

Send and then wait for something, not from same IP bc from\\
inside NAT, but port should be OK
\begin{minted}[fontsize=\small]{python}
pid = os.fork()
if pid:
    print "parent: setting up sniffing"
    # Wait for UDP packet
    data = sniff(filter="udp and port 54040 and net 192.0.2.0/24", count=1)
else:
    time.sleep(10)
    print "child: sending packet"
    sendp(packet,loop=0)
    print "child: closing"
    sys.exit(0)
data[0].show()
\end{minted}

Source port in the inside request packet, becomes the destination port in replies from the server - 54040

\slide{Scapy -- exercise}
% FIXME reuse your Scapy exercise (spoofed source / multi-port)
\begin{list2}
\item \textbf{FIXME:} reuse your Scapy exercise (spoofed source / multi-port)
\end{list2}

\slide{Protection -- Enable More Packet filtering}

\hlkimage{12cm}{network-layers-2022.pdf}

\begin{list2}
\item Packet filtering can be done one single packets -- stateless filtering
\item We can save information about direction and ongoing traffic -- stateful filtering/firewalling
\item \emph{Filtering} can also be setting a maximum number of packets for a protocol -- rate limit by protocol
\end{list2}

\slide{Designing the protection -- bandwidth and rate limit}

\begin{tabularx}{\textwidth-5cm}{|p{5cm}|p{7cm}|X|} \hline
{\bf Protocol} & {\bf Mbps} & {\bf Prefix}\\\hline
TCP & Up to full bandwidth 10Gbps & 192.0.2.0/25 \\\hline
UDP & Less than 1Gbps & 192.0.2.128/25 \\\hline
ICMP & Less than 10Mbps & 192.0.2.0/24 \\\hline
\end{tabularx}


\begin{list2}
\item Create an address plan for your services
\item Monitor your traffic -- how much UDP and TCP do you have, roughly

\item Above is a simplified example -- dig deeper into your traffic
\end{list2}

\slide{Designing the protection -- address families \& protocols }

\begin{tabularx}{\textwidth-5cm}{|p{2cm}|p{2cm}|p{7cm}|X|} \hline
{\bf Address family} & {\bf Protocol} & {\bf Services and ports} & {\bf Prefix}\\\hline
IPv4 & TCP & 25, 80, 8003, 443, 4443 & 192.0.2.0/25 \\\hline
IPv4 & UDP & 53 & 192.0.2.128/25 \\\hline
IPv6 & UDP & 53  & 2001:DB8:ABCD:0053::/60 \\\hline
IPv6 & TCP & 80 443 & 2001:DB8:ABCD:1000::/56 \\\hline
\end{tabularx}

\begin{list2}
\item Direction is also very important -- servers that never initiate connections have fewer requirements
\item How much traffic do you have that uses IPv6 yet? Should an IPv6 DDoS take up all resources?
\item Maybe let IPv4 only use a part, so at least some customers can visit using IPv6?
\item Maybe you can do an allow list for allocated networks, since not all is used yet
\end{list2}

\slide{Stateless firewall filter limit protocols}

\begin{alltt}\footnotesize
term limit-icmp \{
    from \{
        protocol icmp;
    \}
    then \{
        policer ICMP-100M;
        accept;
    \}
\}
term limit-udp \{
    from \{
        protocol udp;
    \}
    then \{
        policer UDP-1000M;
        accept;
    \}
\}
\end{alltt}

Routers also have extensive Class-of-Service (CoS) tools today, and in general rate limiting stuff is nice

\slide{Strict filtering for some servers, still stateless!}

\begin{alltt}\footnotesize
term some-server-allow \{
    from \{
        destination-address \{
            192.0.2.0/25;
        \}
        protocol tcp;
        destination-port [ 25 80 8003 443 4443 ];
    \} then accept;
\}
term some-server-block-unneeded \{
    from \{
        destination-address \{
            192.0.2.0/25; \}
        protocol-except icmp;  \}
    then \{ count some-server-block; discard;
    \}
\}
\end{alltt}

Wut - no UDP, yes only TCP service is used on these servers

\slide{Results from implementing -- DDoS traffic before filtering}
\hlkimage{20cm}{ddos-before-filtering}

\centerline{Only two links shown, at least 3Gbit incoming for this single IP}

\slide{DDoS traffic after filtering}
\hlkimage{20cm}{ddos-after-filtering}

\begin{list1}
\item Link toward server (next level firewall actually) about ~350Mbit outgoing
\item Knowing what it going on, is half the battle
\end{list1}

\slide{Make incremental changes}

\begin{center}
\begin{tikzpicture}[->,>=stealth',scale=0.7, transform shape]
\newlength{\boxwidth}
\setlength{\boxwidth}{0.21\paperwidth}
\newlength{\boxheight}
\setlength{\boxheight}{0.25\paperheight}
\newlength{\boxspace}
\setlength{\boxspace}{10cm}
\input{incremental-changes-flowchart.tikz}
\end{tikzpicture}
\end{center}

\slide{Interesting findings}

\begin{alltt}\small
| snmp-info:
|   enterprise: ciscoSystems
|   engineIDFormat: mac
|   engineIDData: 00:08:4f:xx:yy:zz
|   snmpEngineBoots: 4
|_  snmpEngineTime: {\bf 732d07h09m04s}
\end{alltt}


\begin{list2}
\item Customers have ISPs which have routers with 1800 days uptime\\
Seen using SNMP -- so missing critical updates, and we can affect performance
\item Even big companies still have 1Gbit internet connections, recommend getting 10Gbps
\item Most customers fail at least 3 scenarios which can be countered \emph{without loss of functionality}\\
Example we used a stateless filter to only allow TCP traffic towards a /24 with web services\\
(which don't need or want 10Gbps ICMP/UDP)

\item Ohh we lost our VPN into the environment, ohh the fw console is dead
\end{list2}

\slide{Conclusion DDoS and network attacks}

\hlkimage{10cm}{network-layers-2022.pdf}
~
\begin{list2}
\item You really should try testing, and investigate your existing devices
all of them
\item Choose which devices does which part -- discard early to free resources for later devices to dig deeper
\item This is just one small part of your security posture, extra slides has my take on enterprise network security
\end{list2}

\myquestionspage

\slide{Summary \& references}
\begin{list2}
\item Attacks are easy to generate; defences are filtering, rate limits, monitoring and DDoS protection
\item \textbf{FIXME:} add a CDN / DDoS-protection-service slide if you want it
\end{list2}

\end{document}
